% GENERATED FILE. Edit canonical lessons, then run npm run generate:books.
% canonical-source-hash: fnv1a64:f35ce8438d524e1a
% canonical-lessons: HI-C452-bhavna, HI-C452-ehsas, HI-C452-anubhuti, HI-C452-mahsuskarna, HI-C452-umarna

\chapter{Naming a Feeling}
\label{ch:hi-naming-a-feeling}
\hlchaptermodality{452}

\begin{chapteropening}
\textbf{By the end of this chapter:} \emph{I can name a feeling and say that I felt it.}

Complete the last lesson of chapter 452: I can name a feeling and say that I felt it.
\end{chapteropening}

\section[\texorpdfstring{\hi{भावना} (bhāvnā)}{भावना (bhāvnā)}]{\hi{भावना} (bhāvnā) — a feeling; an emotion}

\label{lesson:HI-C452-bhavna}



\begin{quote}
Before the new one: say the Hindi for turning over (in bed), then the Hindi for a stretch (of the body).
\end{quote}

\subsection*{You'll want to know: \hi{भावना}}
\textbf{\hi{भावना}} — \emph{bhāvnā} — \textquotedblleft{}a feeling; an emotion\textquotedblright{}.

\begin{quote}
\hi{हर} \hi{इंसान} \hi{की} \hi{भावना} \hi{अलग} \hi{होती} \hi{है।}
\end{quote}

\begin{quote}
\emph{har insān kī bhāvnā alag hotī hai.}
\end{quote}

\begin{quote}
Every person's feelings are different.
\end{quote}

Feminine.

\begin{cousinweb}
\textbf{\hi{भावना}} (\emph{bhāvnā}) is a tatsama, borrowed from Sanskrit \emph{bhāvanā}, which Monier-Williams gives with senses such as imagining and a feeling of devotion (Monier-Williams). Hindi uses it for any feeling.
\end{cousinweb}

\begin{grammarlens}[title={What you've built}]
One more word for describing what you lived and felt.
\end{grammarlens}

\subsection*{Your turn}
\begin{itemize}
\raggedright
  \item \emph{Say it:} \emph{bhāvnā}
  \item \emph{Say it:} \emph{bhāvnā}, once more
  \item \emph{Say it:} \emph{har insān kī bhāvnā alag hotī hai.} — Every person's feelings are different.
  \item \emph{Recall:} say the Hindi for turning over (in bed), then the Hindi for a stretch (of the body), then say \emph{bhāvnā} again
\end{itemize}

\begin{tcolorbox}[breakable,colback=teal!4,colframe=teal!35!black,title={Before you move on}]
What does \hi{भावना} mean? (\textquotedblleft{}A feeling; an emotion\textquotedblright{}.) Say it once more.
\end{tcolorbox}
\section[\texorpdfstring{\hi{एहसास} (ehsās)}{एहसास (ehsās)}]{\hi{एहसास} (ehsās) — an awareness; a sense of something}

\label{lesson:HI-C452-ehsas}



\begin{quote}
Before the new one: say the Hindi for a stretch (of the body), then the Hindi for a feeling.
\end{quote}

\subsection*{You'll want to know: \hi{एहसास}}
\textbf{\hi{एहसास}} — \emph{ehsās} — \textquotedblleft{}an awareness; a sense of something\textquotedblright{}.

\begin{quote}
\hi{अरुण} \hi{को} \hi{अपनी} \hi{ग़लती} \hi{का} \hi{एहसास} \hi{हुआ।}
\end{quote}

\begin{quote}
\emph{Aruṇ ko apnī ġaltī kā ehsās huā.}
\end{quote}

\begin{quote}
Arun realized he had made a mistake.
\end{quote}

Masculine.

\begin{cousinweb}
\textbf{\hi{एहसास}} (\emph{ehsās}) is usually traced to Arabic \emph{iḥsās}, \textquotedblleft{}perceiving, feeling\textquotedblright{}, which reached Hindi through Persian.
\end{cousinweb}

\begin{grammarlens}[title={What you've built}]
One more word for describing what you lived and felt.
\end{grammarlens}

\subsection*{Your turn}
\begin{itemize}
\raggedright
  \item \emph{Say it:} \emph{ehsās}
  \item \emph{Say it:} \emph{ehsās}, once more
  \item \emph{Say it:} \emph{Aruṇ ko apnī ġaltī kā ehsās huā.} — Arun realized he had made a mistake.
  \item \emph{Recall:} say the Hindi for a stretch (of the body), then the Hindi for a feeling, then say \emph{ehsās} again
\end{itemize}

\begin{tcolorbox}[breakable,colback=teal!4,colframe=teal!35!black,title={Before you move on}]
What does \hi{एहसास} mean? (\textquotedblleft{}An awareness; a sense of something\textquotedblright{}.) Say it once more.
\end{tcolorbox}
\section[\texorpdfstring{\hi{अनुभूति} (anubhūti)}{अनुभूति (anubhūti)}]{\hi{अनुभूति} (anubhūti) — a deep feeling; an inner sense}

\label{lesson:HI-C452-anubhuti}



\begin{quote}
Before the new one: say the Hindi for a feeling, then the Hindi for an awareness.
\end{quote}

\subsection*{You'll want to know: \hi{अनुभूति}}
\textbf{\hi{अनुभूति}} — \emph{anubhūti} — \textquotedblleft{}a deep feeling; an inner sense\textquotedblright{}.

\begin{quote}
\hi{गाना} \hi{सुनकर} \hi{मीरा} \hi{को} \hi{शांति} \hi{की} \hi{अनुभूति} \hi{हुई।}
\end{quote}

\begin{quote}
\emph{gānā sunkar Mīrā ko śānti kī anubhūti huī.}
\end{quote}

\begin{quote}
Listening to the song, Mira felt a deep sense of peace.
\end{quote}

Feminine.

\begin{cousinweb}
\textbf{\hi{अनुभूति}} (\emph{anubhūti}) is a tatsama, borrowed from Sanskrit \emph{anubhūti}, \textquotedblleft{}perception, knowledge\textquotedblright{} (Monier-Williams). It is a formal word for a feeling sensed deep inside; \emph{anubhav} is the experience itself.
\end{cousinweb}

\begin{grammarlens}[title={What you've built}]
One more word for describing what you lived and felt.
\end{grammarlens}

\subsection*{Your turn}
\begin{itemize}
\raggedright
  \item \emph{Say it:} \emph{anubhūti}
  \item \emph{Say it:} \emph{anubhūti}, once more
  \item \emph{Say it:} \emph{gānā sunkar Mīrā ko śānti kī anubhūti huī.} — Listening to the song, Mira felt a deep sense of peace.
  \item \emph{Recall:} say the Hindi for a feeling, then the Hindi for an awareness, then say \emph{anubhūti} again
\end{itemize}

\begin{tcolorbox}[breakable,colback=teal!4,colframe=teal!35!black,title={Before you move on}]
What does \hi{अनुभूति} mean? (\textquotedblleft{}A deep feeling; an inner sense\textquotedblright{}.) Say it once more.
\end{tcolorbox}
\section[\texorpdfstring{\hi{महसूस} \hi{करना} (mahsūs karnā)}{महसूस करना (mahsūs karnā)}]{\hi{महसूस} \hi{करना} (mahsūs karnā) — to feel}

\label{lesson:HI-C452-mahsuskarna}



\begin{quote}
Before the new one: say the Hindi for an awareness, then the Hindi for a deep feeling.
\end{quote}

\subsection*{You'll want to know: \hi{महसूस} \hi{करना}}
\textbf{\hi{महसूस} \hi{करना}} — \emph{mahsūs karnā} — \textquotedblleft{}to feel\textquotedblright{}.

\begin{quote}
\hi{मीरा} \hi{ने} \hi{कंधे} \hi{में} \hi{दर्द} \hi{महसूस} \hi{किया।}
\end{quote}

\begin{quote}
\emph{Mīrā ne kandhe meṁ dard mahsūs kiyā.}
\end{quote}

\begin{quote}
Mira felt a pain in her shoulder.
\end{quote}

Past, with \hi{ने} (he or she …): \hi{उसने} \hi{दर्द} \hi{महसूस} \hi{किया}, \emph{usne dard mahsūs kiyā}.

\begin{cousinweb}
\textbf{\hi{महसूस} \hi{करना}} (\emph{mahsūs karnā}) joins \emph{mahsūs}, usually traced to Arabic \emph{maḥsūs}, \textquotedblleft{}perceived, felt\textquotedblright{}, which came through Persian, with \emph{karnā}, \textquotedblleft{}to do\textquotedblright{}: to do the feeling of something. The thing felt is the object, so the perfective agrees with it.
\end{cousinweb}

\begin{grammarlens}[title={What you've built}]
One more word for describing what you lived and felt.
\end{grammarlens}

\subsection*{Your turn}
\begin{itemize}
\raggedright
  \item \emph{Say it:} \emph{mahsūs karnā}
  \item \emph{Say it:} \emph{mahsūs karnā}, once more
  \item \emph{Say it:} \emph{Mīrā ne kandhe meṁ dard mahsūs kiyā.} — Mira felt a pain in her shoulder.
  \item \emph{Recall:} say the Hindi for an awareness, then the Hindi for a deep feeling, then say \emph{mahsūs karnā} again
\end{itemize}

\begin{tcolorbox}[breakable,colback=teal!4,colframe=teal!35!black,title={Before you move on}]
What does \hi{महसूस} \hi{करना} mean? (\textquotedblleft{}To feel\textquotedblright{}.) Say it once more.
\end{tcolorbox}
\section[\texorpdfstring{\hi{उमड़ना} (umaṛnā)}{उमड़ना (umaṛnā)}]{\hi{उमड़ना} (umaṛnā) — to well up; to surge}

\label{lesson:HI-C452-umarna}



\begin{quote}
Before the new one: say the Hindi for a deep feeling, then the Hindi for to feel.
\end{quote}

\subsection*{You'll want to know: \hi{उमड़ना}}
\textbf{\hi{उमड़ना}} — \emph{umaṛnā} — \textquotedblleft{}to well up; to surge\textquotedblright{}.

\begin{quote}
\hi{बेटी} \hi{को} \hi{देखकर} \hi{माँ} \hi{की} \hi{आँखों} \hi{में} \hi{आँसू} \hi{उमड़} \hi{आए।}
\end{quote}

\begin{quote}
\emph{beṭī ko dekhkar māṁ kī ā̃khoṁ meṁ ā̃sū umaṛ āe.}
\end{quote}

\begin{quote}
When she saw her daughter, tears welled up in the mother's eyes.
\end{quote}

Past (it …): \hi{वह} \hi{नदी} \hi{उमड़} \hi{पड़ी}, \emph{vah nadī umaṛ paṛī}.

\begin{cousinweb}
\textbf{\hi{उमड़ना}} (\emph{umaṛnā}) is a tadbhava: CDIAL lists Hindi \emph{umaḍnā}, \textquotedblleft{}to swell, heave\textquotedblright{}, under \emph{*ummaḍ}, a form it reconstructs from its descendants (CDIAL 2344). Feelings, tears, rivers and crowds all well up with it.
\end{cousinweb}

\begin{grammarlens}[title={What you've built}]
One more word for describing what you lived and felt.
\end{grammarlens}

\subsection*{Your turn}
\begin{itemize}
\raggedright
  \item \emph{Say it:} \emph{umaṛnā}
  \item \emph{Say it:} \emph{umaṛnā}, once more
  \item \emph{Say it:} \emph{beṭī ko dekhkar māṁ kī ā̃khoṁ meṁ ā̃sū umaṛ āe.} — When she saw her daughter, tears welled up in the mother's eyes.
  \item \emph{Recall:} say the Hindi for a deep feeling, then the Hindi for to feel, then say \emph{umaṛnā} again
\end{itemize}

\begin{tcolorbox}[breakable,colback=teal!4,colframe=teal!35!black,title={Before you move on}]
What does \hi{उमड़ना} mean? (\textquotedblleft{}To well up; to surge\textquotedblright{}.) Say it once more.
\end{tcolorbox}
